\section{厄米多项式与拉盖尔多项式}
\idx{厄米多项式}{emi}
\idx{拉盖尔多项式}{lageer}
球坐标里的角向方程给出勒让德多项式和球谐函数,那一部分已经在曲面坐标一章里.中心力场的径向方程,以及简谐振子,还会再产生两类多项式.它们都是在无穷区间上带权重正交的 Sturm--Liouville 本征函数.

\subsection{量子谐振子}
一维谐振子
\begin{equation}
  -\frac{\hbar^2}{2m}\frac{\mathrm{d}^2\psi}{\mathrm{d}x^2}
  +\half m\omega^2 x^2\,\psi
  =E\psi.
\end{equation}
用 $\xi=x\sqrt{m\omega/\hbar}$ 和 $\varepsilon=2E/(\hbar\omega)$ 把方程写成
\begin{equation}
  -\frac{\mathrm{d}^2\psi}{\mathrm{d}\xi^2}+(\xi^2-\varepsilon)\psi=0.
\end{equation}
$\xi\to\infty$ 时势能项主导,可积的解必须像 $e^{-\xi^2/2}$ 那样衰减.令 $\psi(\xi)=v(\xi)\,e^{-\xi^2/2}$,代入后 $v$ 满足
\begin{equation}
  v''-2\xi v'+(\varepsilon-1)v=0.
  \label{eq:hermite-ode}
\end{equation}
在 $\xi=0$ 处把 $v$ 写成幂级数.递推每一步把次数提高 $2$,一般项像 $\xi^{2k}/(k!)$ 再乘一个不趋于零的因子,这样 $v$ 会增长得比 $e^{\xi^2/2}$ 还快,$\psi$ 在无穷远处不可积.级数必须在有限项截断.截断条件是
\begin{equation}
  \varepsilon-1=2n, \qquad n=0,1,2,\ldots,
\end{equation}
因此能级
\begin{equation}
  E_n=\hbar\omega\bigl(n+\half\bigr).
  \label{eq:oscillator-energy}
\end{equation}
截断后的多项式就是厄米多项式.罗德里格斯公式
\begin{equation}
  H_n(\xi)=(-1)^n e^{\xi^2}\frac{\mathrm{d}^n}{\mathrm{d}\xi^n}e^{-\xi^2}
\end{equation}
给出 $H_0=1$,$H_1=2\xi$,$H_2=4\xi^2-2$,$H_3=8\xi^3-12\xi$.它满足
\begin{equation}
  H_n''-2\xi H_n'+2n H_n=0,
\end{equation}
正是\eqref{eq:hermite-ode}在 $\varepsilon-1=2n$ 时的方程.生成函数
\begin{equation}
  e^{2\xi t-t^2}=\sum_{n=0}^{\infty} H_n(\xi)\frac{t^n}{n!}
\end{equation}
对 $t$ 求导可以得到递推 $H_{n+1}=2\xi H_n-2n H_{n-1}$.

权重 $e^{-\xi^2}$ 下它们正交,
\begin{equation}
  \int_{-\infty}^{\infty} e^{-\xi^2} H_m(\xi) H_n(\xi)\,\mathrm{d}\xi
  =\sqrt{\pi}\, 2^n n!\, \delta_{mn}.
\end{equation}
归一化的束缚态是
\begin{equation}
  \psi_n(x)
  =\bigl(2^n n!\bigr)^{-1/2}
    \bigl(\tfrac{m\omega}{\pi\hbar}\bigr)^{1/4}
    H_n(\xi)\, e^{-\xi^2/2},
\end{equation}
其中 $\xi=x\sqrt{m\omega/\hbar}$.含时解再乘 $e^{-\imath E_n t/\hbar}$.基态 $n=0$ 没有节点,$H_0=1$,概率密度是一个高斯.

\begin{example}
用罗德里格斯公式算 $H_2$,并验证它满足厄米方程.
\end{example}
\begin{solution}
\begin{equation}
  \frac{\mathrm{d}}{\mathrm{d}\xi}e^{-\xi^2}=-2\xi e^{-\xi^2},
  \qquad
  \frac{\mathrm{d}^2}{\mathrm{d}\xi^2}e^{-\xi^2}=(4\xi^2-2)e^{-\xi^2}.
\end{equation}
因此 $H_2=(-1)^2 e^{\xi^2}(4\xi^2-2)e^{-\xi^2}=4\xi^2-2$.
\begin{equation}
  H_2'=8\xi, \qquad H_2''=8,
  \qquad
  H_2''-2\xi H_2'+4 H_2=8-2\xi\cdot 8\xi+4(4\xi^2-2)=0.
\end{equation}
\end{solution}

\subsection{广义拉盖尔多项式}
氢原子的库仑场是中心力.波函数写成 $R(r)Y_{\ell m}(\theta,\phi)$ 之后,角向部分就是球谐函数.径向函数 $u(r)=r R(r)$ 满足
\begin{equation}
  -\frac{\hbar^2}{2m}u''
  +\Bigl(-\frac{e^2}{r}+\frac{\hbar^2\ell(\ell+1)}{2m r^2}\Bigr)u
  =E u.
\end{equation}
束缚态 $E<0$.令
\begin{equation}
  \kappa=\sqrt{-2mE}/\hbar,
  \qquad
  \rho=2\kappa r,
  \qquad
  n=\frac{m e^2}{\hbar^2\kappa},
\end{equation}
这里的 $n$ 先只是一个参数.在 $\rho\to\infty$ 时 $u\sim e^{-\rho/2}$,在原点 $u\sim\rho^{\ell+1}$.把这两个因子提出去,
\begin{equation}
  u(r)=\rho^{\ell+1} e^{-\rho/2} w(\rho),
\end{equation}
$w$ 满足广义拉盖尔方程.可积性再次要求级数截断,参数 $n$ 必须是正整数,并且 $n\geq\ell+1$.能级因此是
\begin{equation}
  E_n=-\frac{m e^4}{2\hbar^2}\frac{1}{n^2}.
\end{equation}
用高斯单位制,数值上大约是 $-13.6\,\mathrm{eV}/n^2$.能量只依赖 $n$,不依赖 $\ell$.对固定的 $n$,角动量可以取 $\ell=0,1,\ldots,n-1$,每个 $\ell$ 又有 $2\ell+1$ 个 $m$.简并度是
\begin{equation}
  \sum_{\ell=0}^{n-1}(2\ell+1)=n^2.
\end{equation}
这是库仑势 $1/r$ 特有的简并.微扰一般会把这 $n^2$ 个态拆开.

截断后的多项式是广义拉盖尔多项式.罗德里格斯公式
\begin{equation}
  L_n^{(\alpha)}(x)
  =\frac{e^{x} x^{-\alpha}}{n!}
    \frac{\mathrm{d}^n}{\mathrm{d}x^n}\bigl(e^{-x} x^{n+\alpha}\bigr),
  \qquad \alpha>-1.
\end{equation}
$\alpha=0$ 时就是普通拉盖尔多项式 $L_n$.前几个是
\begin{equation}
  L_0^{(\alpha)}=1,
  \qquad
  L_1^{(\alpha)}=1+\alpha-x,
  \qquad
  L_2^{(\alpha)}=\half\bigl[(1+\alpha)(2+\alpha)-2(2+\alpha)x+x^2\bigr].
\end{equation}
正交关系是
\begin{equation}
  \int_0^{\infty} e^{-x} x^{\alpha} L_n^{(\alpha)}(x) L_m^{(\alpha)}(x)\,\mathrm{d}x
  =\frac{\Gamma(n+\alpha+1)}{n!}\,\delta_{mn}.
\end{equation}
氢原子径向波函数用到的是
\begin{equation}
  w(\rho)=L_{n-\ell-1}^{(2\ell+1)}(\rho),
  \qquad
  \rho=\frac{2r}{n a_0},
  \qquad
  a_0=\frac{\hbar^2}{m e^2}.
\end{equation}
$n-\ell-1$ 是节点数.基态 $n=1,\ell=0$ 没有节点,$L_0^{(1)}=1$,径向概率密度 $r^2|R|^2$ 在 $r=a_0$ 取最大.
